% Source: https://artofproblemsolving.com/community/c3h1527002_factoring_quadratics_help
% Stable identity: aops-topic-1527002-q01
If a product of two numbers is $0$, at least one of the factors is $0$ (the zero-product property). So $(x-8)(x+5)=0$ means
\[
x-8=0\quad\text{or}\quad x+5=0,\qquad\text{i.e.}\qquad x=8\quad\text{or}\quad x=-5 .
\]
The sign ``flips'' because the solution makes the factor zero: $x+5=0$ needs $x=-5$, not $x=5$. Check: $8^2-3\cdot8-40=64-24-40=0$ and $(-5)^2-3(-5)-40=25+15-40=0$. The solutions are $\boxed{x=8\text{ and }x=-5}$.
